% Generated draft: requires mathematical and source review.
\documentclass[11pt]{article}
\usepackage[margin=25mm]{geometry}
\usepackage{fontspec}
\setmainfont{Latin Modern Roman}
\usepackage{amsmath,amssymb,mathtools,unicode-math}
\setmathfont{Latin Modern Math}
\usepackage{enumitem,needspace,xurl,hyperref}
\setlength{\parindent}{0pt}
\setlength{\parskip}{5pt}
\setlength{\emergencystretch}{3em}
\newcommand{\N}{\mathbb{N}}
\newcommand{\Z}{\mathbb{Z}}
\newcommand{\Q}{\mathbb{Q}}
\newcommand{\R}{\mathbb{R}}
\newcommand{\C}{\mathbb{C}}
\newcommand{\F}{\mathbb{F}}
\newcommand{\A}{\mathbb{A}}
\newcommand{\PP}{\mathbb{P}}
\newcommand{\E}{\mathbb{E}}
\newcommand{\eps}{\varepsilon}
\newcommand{\dd}{\,\mathrm{d}}
\begin{document}
\section*{UnsolvedMath ID: 11000212}
\textbf{Source problem code:} catalogue entry
\subsection{Definitions and mathematical statement}
A Fuchsian group is a discrete subgroup of $\operatorname{PSL}_2(\R)$. In the usual flat-surface setting, a nonzero quadratic differential $q$ on a Riemann surface determines a Teichm\"uller disc $D_q\subset\operatorname{Teich}(S)$. Its setwise stabilizer is
\[
 \operatorname{Stab}_{\operatorname{Mod}(S)}(D_q)
 =\{\varphi\in\operatorname{Mod}(S):\varphi(D_q)=D_q\}.
\]
The induced action on the disc gives a projective derivative group in $\operatorname{PSL}_2(\R)$, customarily called the Veech group. In a translation-surface convention the derivative group is initially in $\operatorname{SL}_2(\R)$. The disc stabilizer can have a finite kernel in its action on the disc, so it need not be abstractly identical to the projectivized derivative group.

Problem 5 asks two realization questions: which Fuchsian groups occur as Veech groups, and which subgroups of the mapping class group occur as Veech groups, equivalently as stabilizers of Teichm\"uller discs in the source convention? The intended object is the full Veech group or full stabilizer, not merely a subgroup contained in one.

The excerpt does not specify genus, surface type, differential stratum, translation versus half-translation surfaces, or compactness/finite-area conventions. It also does not say whether realization is up to abstract isomorphism or conjugacy in $\operatorname{PSL}_2(\R)$, and analogously in the mapping class group. No finite-covolume or lattice restriction is stated. The source's equivalence between the two questions must be understood with its derivative/projectivization and finite-kernel conventions fixed.
\subsection{Short English statement}
Classify the Fuchsian groups arising as Veech groups and the mapping-class subgroups that are full Teichm\"uller-disc stabilizers, keeping track of derivative and finite-kernel conventions.
\subsection{Sources}
\begin{itemize}
\item Local registry: \url{data/unsolvedmath/problems.json}, permanent UnsolvedMath ID $11000212$, exact source code \texttt{catalogue entry}.
\item Farb, \emph{Problems on Mapping Class Groups and Related Topics}, Problem 5, PDF page 259: \url{https://www.math.uchicago.edu/~farb/papers/mcgbook.pdf}. Supplied access date: 2026-07-29; extraction: pdf-font-structure; repeated local numbering disambiguated by page among 329 headings.
\item No additional primary citation or individual attribution is supplied. The registry explicitly cannot verify a settled literature answer.
\end{itemize}
\end{document}
